\Cref{part:constrsync} deals with the constrained synchronization problem and is based on the works~\cite{FernauGHHVW19,
	DBLP:conf/cocoon/Hoffmann20,
	DBLP:conf/ictcs/Hoffman20,
	DBLP:conf/fct/Hoffmann21,
	DBLP:journals/tcs/Hoffmann21,
	DBLP:conf/cocoon/Hoffmann21,
	DBLP:conf/wia/Hoffmann22,
	HoffmannCocoonExtended}.
The results from~\cite{DBLP:conf/cocoon/Hoffmann20,DBLP:conf/fct/Hoffmann21,DBLP:journals/tcs/Hoffmann21,DBLP:conf/cocoon/Hoffmann21,HoffmannCocoonExtended}
were revisited with an emphasis on a more streamlined, coherent, and simplified presentation.
This was achieved primarily through the use of general methods introduced in~\cite{DBLP:conf/cocoon/Hoffmann21,HoffmannCocoonExtended}.

In \cref{constr_sync:sec:motivation}, we give motivational examples for the constrained synchronization problem.

\Cref{constr_sync:chap:further_notions} introduces additional notions that we will need in \cref{part:constrsync}. We also present a new characterization of solid codes that, to the best of our knowledge, has not been published before.

\Cref{constr_sync:chap:gen_results} presents general results on manipulating constraint languages while preserving
the complexity or to deduce complexity lower bounds. Furthermore, we give sufficient conditions for \NP- and \PSPACE-hardness. 
Some of these results
were presented in~\cite{DBLP:conf/cocoon/Hoffmann21,HoffmannCocoonExtended},
which generalized results from~\cite{FernauGHHVW19},
and some results are new.

The central question investigated in \cref{constr_sync:chap:classes_constraint_languages}
is the computational complexity of \constrsync{L} for different classes of regular
constraint languages $L \subseteq \Sigma^*$.
While a complete classification remains an \emph{open problem},
we present several partial results by giving complete classifications for the following classes of constraint languages:
letter-bounded languages, languages accepted by partial automata (PDFAs, introduced in
\cref{constr_sync:chap:further_notions}) with at most two states over an arbitrary alphabet and at most three states over a binary alphabet, and commutative languages.
We also obtain results for infinite, sparse, and monoid languages, as well as
for languages accepted by strongly connected PDFAs.
The classification for commutative constraint languages appeared
in~\cite{DBLP:conf/cocoon/Hoffmann20,DBLP:journals/tcs/Hoffmann21}.
Here, we present a simplified formulation and proof.
The results for sparse and letter-bounded languages appeared
in~\cite{DBLP:conf/ictcs/Hoffman20,DBLP:conf/fct/Hoffmann21}.

Lastly, motivated by the question of what features are needed in input automata to realize certain computational complexities, in \cref{constr_sync:chap:constr_input} we investigate specific types of input automata. Some of the results from this chapter appeared in~\cite{DBLP:conf/wia/Hoffmann22}, and some mistakes from this paper have been corrected.